% WHAT RANGE A BREAK PENALTY IS BROUGHT INTO.
%
% The penalty at a breakpoint is brought into range before anything is done
% with it. Ten thousand or more forbids the break outright; ten thousand or
% more the other way IS exactly -10000 from there on, whatever the document
% wrote. So
%
%   A\penalty-1000000000 A\par
%
% traces `p=-10000', not `p=-1000000000' -- and, because -10000 is what the
% forced break is recognised BY, the break is taken as a forced one and its
% demerits are the artificial `d=*' rather than a number.
%
% HSTeX carried the penalty as written, so a penalty below the forced one
% was traced as itself and weighed as an ordinary penalty. trip line 147
% writes \penalty-1000000000 in a paragraph.
\catcode`\{=1 \catcode`\}=2
\tracingonline=1 \tracingparagraphs=1
\font\rip=trip \rip
\hsize=100pt \parindent=0pt \hbadness=10000
\message{[A a penalty far below the forced one]}
\setbox0=\vbox{A\penalty-1000000000 A\par}
\message{[B the forced one itself]}
\setbox0=\vbox{A\penalty-10000 A\par}
\message{[C a penalty that forbids a break]}
\setbox0=\vbox{A\penalty10000 A\par}
\message{[D one far above it]}
\setbox0=\vbox{A\penalty1000000000 A\par}
\message{[E an ordinary one]}
\setbox0=\vbox{A\penalty50 A\par}
\message{[done]}
\end
